% Authority: German Wikiversity pageid 152873, revid 1094182, 2026-06-14T15:10:42Z.
% Exact UTF-8 wikitext witness SHA-256: 1050f074ecb58b7a8211d756b7fcd6fe4a0a4db6a024083c31ca3c5469e9fd50.
% Expanded German TeX source SHA-256: 6896eb6a4a6f4c25bcfaab674847dcd78bc079cfd6c58b3b8542c5c220d85e7b.

Untuk memperoleh
\definitionsverweis {medan normal satuan}{}{,}
kita mulai dari gradien fungsi yang mendeskripsikan elips tersebut, yaitu

\mavergleichskettedisp
{\vergleichskette
{ G(x,y)
}
{ =} { \begin{pmatrix}  2 \alpha x  \\ 2 \beta y    \end{pmatrix}
}
{ } {
}
{ } {
}
{ } {
}
}
{}{}{}
dan dengan menormalisasikannya kita memperoleh

\mavergleichskettedisp
{\vergleichskette
{ N(x,y)
}
{ =} {  { \frac{   1 }{  \sqrt{ \alpha^2 x^2 + \beta^2 y^2 }  } } \begin{pmatrix}   \alpha x  \\  \beta y    \end{pmatrix}
}
{ } {
}
{ } {
}
{ } {
}
}
{}{}{.}
Dengan demikian, pemetaan Gauss tersebut adalah
\maabbeledisp {\varphi} {Y = \left\{  \begin{pmatrix}  x  \\ y   \end{pmatrix}   \in \R^2  \mid   \alpha x^2+\beta y^2 = 1        \right\}  } { S^1
} { \begin{pmatrix}  x  \\ y   \end{pmatrix} } { { \frac{   1 }{  \sqrt{ \alpha^2 x^2 + \beta^2 y^2 }  } } \begin{pmatrix}   \alpha x  \\  \beta y    \end{pmatrix}
} {.} Kesurjektifan pemetaan ini mengikuti
[[Hyperfläche/Glatt und kompakt/Gauß-Abbildung surjektiv/Fakt|Teorema 2.11]].
Untuk membuktikan injektivitasnya, kita tinjau syarat

\mavergleichskettedisp
{\vergleichskette
{ { \frac{   1  }{  \sqrt{ \alpha^2 x^2 + \beta^2 y^2 }  } } \begin{pmatrix}   \alpha x  \\  \beta y    \end{pmatrix}
}
{ =} { { \frac{   1  }{  \sqrt{ \alpha^2 z^2 + \beta^2 w^2 }  } } \begin{pmatrix}   \alpha z  \\  \beta w    \end{pmatrix}
}
{ } {
}
{ } {
}
{ } {
}
}
{}{}{.}
Ini berarti bahwa
\mathkor {} {\begin{pmatrix}   \alpha x  \\  \beta y    \end{pmatrix}} {dan} {\begin{pmatrix}   \alpha z  \\  \beta w    \end{pmatrix}} {}
bergantung secara linear dengan faktor positif $c$, sehingga
$ (x,y)=c(z,w) $. Karena kedua titik terletak pada $Y$,

\[
1=\alpha x^2+\beta y^2
=c^2\left(\alpha z^2+\beta w^2\right)=c^2.
\]

Karena $c>0$, diperoleh $c=1$. Jadi berlaku

\mavergleichskettedisp
{\vergleichskette
{ \begin{pmatrix}   \alpha x  \\  \beta y    \end{pmatrix}
}
{ =} { \begin{pmatrix}   \alpha z  \\  \beta w    \end{pmatrix}
}
{ } {
}
{ } {
}
{ } {
}
}
{}{}{}
dan oleh karena itu

\mavergleichskette
{\vergleichskette
{ x
}
{ = }{ z
}
{  }{
}
{    }{
}
{   }{
}
}
{}{}{}
dan

\mavergleichskette
{\vergleichskette
{ y
}
{ = }{ w
}
{  }{
}
{    }{
}
{   }{
}
}
{}{}{.}

 [[Kategorie:Latexseite]]
